\section{Chuỗi descriptor và scratchpad ping-pong}\label{sec:5}

Kiến trúc chỉ có một chế độ. CPU dựng danh sách descriptor trong DRAM và ra một lệnh START cho cả
chuỗi. Mỗi buffer A, B, C, BIAS có hai bank. Nạp, tính và ghi của các descriptor kế nhau chạy chồng
trên hai bank.

\subsection{Nguyên tắc}\label{sec:5.1}

\paragraph{Nguyên tắc}\label{par:5-1}
\begin{itemize}
\item Không thêm, bớt hay đổi port top-level (mục~\ref{sec:1.7}). Offset thanh ghi giữ nguyên; 0x08 … 0x24 dự trữ vì tham số của job nằm trong descriptor
\item DMA có kênh đọc và kênh ghi độc lập, có lệnh FETCH riêng. SRAM có bit chọn bank cho từng buffer. Main FSM lập lịch ba engine
\item Tham số của mỗi job lấy từ descriptor. \sig{DESC_PTR} là bản chụp duy nhất lúc START
\end{itemize}

\subsection{DMA theo chuỗi descriptor}\label{sec:5.2}

CPU dựng sẵn một danh sách descriptor trong DRAM, mỗi descriptor mô tả một phép GEMM. CPU ghi địa chỉ
descriptor đầu vào \sig{DESC_PTR}, thực hiện \sig{fence rw,rw}, ghi \sig{CTRL[2]} = 1 (START), rồi
chờ DONE một lần cho cả chuỗi. Accelerator tự đọc từng descriptor và chạy hết chuỗi.

\begin{spectable}{|L{0.1}|L{0.08}|L{0.17}|L{0.12}|L{0.53}|}{Bố cục descriptor (16 word, 64 byte)}{tab:5-2}{\textbf{Offset} & \textbf{Word} & \textbf{Tên} & \textbf{Bit} & \textbf{Nội dung}}
0x00 & 0 & \sig{DIM} & [31:0] & M[7:0], N[15:8], K[23:16]; [31:24] phải bằng 0 \\\hline

0x04 & 1 & \sig{A_ADDR} & [31:0] & Địa chỉ byte của A[0][0] \\\hline

0x08 & 2 & \sig{B_ADDR} & [31:0] & Địa chỉ byte của B[0][0]; bỏ qua khi \sig{REUSE_B} = 1 \\\hline

0x0C & 3 & \sig{C_ADDR} & [31:0] & Địa chỉ byte của C[0][0] hoặc Y[0][0] \\\hline

0x10 & 4 & \sig{LDA}, \sig{LDB} & [15:0], [31:16] & Stride theo hàng của A và B (byte); 0 = các hàng liền nhau \\\hline

0x14 & 5 & \sig{LDC} & [15:0] & Stride của C hoặc Y (byte); 0 = liền nhau; [31:16] dự trữ \\\hline

0x18 & 6 & \sig{FLAGS} & [4:0] & [0] \sig{LAST}, [1] \sig{WAIT_PREV}, [2] \sig{REUSE_B}, [3] \sig{POST_EN}, [4] \sig{RELU_EN}; [31:5] dự trữ \\\hline

0x1C & 7 & \sig{NEXT_PTR} & [31:0] & Địa chỉ descriptor kế tiếp, căn 64 byte; bỏ qua khi \sig{LAST} = 1 \\\hline

0x20 & 8 & \sig{BIAS_ADDR} & [31:0] & Địa chỉ byte của bias[0]; chỉ dùng khi \sig{POST_EN} = 1 \\\hline

0x24 & 9 & \sig{QM0} & [31:0] & Hệ số nhân $M_0$, INT32 có dấu; chỉ dùng khi \sig{POST_EN} = 1 \\\hline

0x28 & 10 & \sig{QSHIFT} & [4:0] & Lượng dịch $sh$, 0 … 31; [31:5] dự trữ \\\hline

0x2C & 11 & \sig{TAG} & [31:0] & Tùy phần mềm; phần cứng không đọc, mọi giá trị hợp lệ \\\hline

0x30–0x3C & 12–15 & — & — & Dự trữ, phải bằng 0 \\\hline

\end{spectable}

Với stride, địa chỉ phần tử là: A[i][k] tại \sig{A_ADDR} + i·\sig{LDA} + k; B[k][j] tại
\sig{B_ADDR} + k·\sig{LDB} + j; C[i][j] tại \sig{C_ADDR} + i·\sig{LDC} + 4j; Y[i][j] tại
\sig{C_ADDR} + i·\sig{LDC} + j. Stride = 0 được hiểu là đúng độ dài hàng.

\begin{notebox}[Nguồn]
Bố cục 16 word ở bảng~\ref{tab:5-2} chỉ có trong \sig{spec_1}. \sig{spec_2} xác nhận descriptor dài
16 word và 64 byte, xác nhận lệnh FETCH (\sig{rows} = 1, \sig{rwords} = 16, \sig{stride} = 64),
xác nhận yêu cầu căn 64 byte và xác nhận các cờ \sig{LAST}, \sig{WAIT_PREV}, \sig{REUSE_B} cùng
word \sig{TAG} bị bỏ qua, nhưng \textbf{không} nêu vị trí từng word và từng bit. Vì vậy baseline lấy
nguyên bố cục của \sig{spec_1}; không có mâu thuẫn cần giải, chỉ có một nguồn chi tiết hơn.
\end{notebox}

\subsubsection{Lấy và tiêu thụ descriptor}\label{sec:5.2.1}

\begin{spectable}{|L{0.22}|L{0.78}|}{Lấy và tiêu thụ một descriptor}{tab:5-3}{\textbf{Bước} & \textbf{Nội dung}}
Phát lệnh FETCH & Main FSM phát lệnh kênh đọc với \sig{o_rd_dst} = 3, \sig{o_rd_base} = con trỏ descriptor, \sig{o_rd_rows} = 1, \sig{o_rd_rwords} = 16, \sig{o_rd_stride} = 64, \sig{o_rd_bank} = 0 \\\hline

Truyền AXI & Một burst INCR 16 beat, \sig{m_axi_arlen} = 15, \sig{m_axi_arsize} = 3'b010. Burst này không bao giờ vượt biên 4 KiB vì địa chỉ căn 64 byte và độ dài 64 byte \\\hline

Trả về & Descriptor \textbf{không} đi vào scratchpad. DMA trả từng word qua \sig{o_desc_valid}, \sig{o_desc_idx[3:0]} = 0 … 15, \sig{o_desc_data[31:0]} \\\hline

Nhận & Main FSM ghi 16 word vào bộ đệm descriptor nội bộ. Xong khi nhận đủ word thứ 15, hoặc lỗi mã 4 \\\hline

Kiểm tra & Một chu kỳ: \sig{gemm_job_check} dựng \sig{job_t} và kiểm điều kiện hợp lệ ở mục~\ref{sec:5.2.2}. Sai → \sig{ERR_CODE} = 1 \\\hline

Con trỏ kế tiếp & Lấy từ \sig{NEXT_PTR} của descriptor vừa kiểm. \sig{LAST} = 1 kết thúc chuỗi, \sig{NEXT_PTR} bị bỏ qua \\\hline

Thứ tự & Mỗi descriptor được đọc \textbf{đúng một lần}, theo đúng thứ tự chuỗi. Engine Đọc làm tuần tự FETCH(d) → kiểm tra(d) → mọi LOAD của d → FETCH(d+1) \\\hline

\end{spectable}

\subsubsection{Điều kiện descriptor hợp lệ}\label{sec:5.2.2}

Descriptor hợp lệ khi \textbf{tất cả} các điều kiện sau đúng. Sai bất kỳ điều kiện →
\sig{ERR_CODE} = 1 và descriptor đó là descriptor lỗi.

\begin{spectable}{|L{0.26}|L{0.74}|}{Điều kiện descriptor hợp lệ}{tab:5-4}{\textbf{Hạng mục} & \textbf{Điều kiện}}
Kích thước & 1 ≤ M ≤ 32; N, K ∈ \{4, 8, …, 32\} \\\hline

Bit dự trữ của \sig{DIM} & \sig{DIM[31:24]} = 0 \\\hline

Căn chỉnh địa chỉ & Mọi địa chỉ \textbf{được dùng} căn 4 byte: \sig{A_ADDR}, \sig{C_ADDR}, \sig{B_ADDR} (trừ khi \sig{REUSE_B} = 1), \sig{BIAS_ADDR} (chỉ khi \sig{POST_EN} = 1) \\\hline

Con trỏ kế tiếp & \sig{NEXT_PTR} căn 64 byte khi \sig{LAST} = 0 \\\hline

Stride & Mỗi stride (\sig{LDA}, \sig{LDB}, \sig{LDC}) bằng 0, hoặc đồng thời ≥ độ dài hàng và là bội của 4. Độ dài hàng: K cho A, N cho B, 4N cho C khi \sig{POST_EN} = 0, N cho Y khi \sig{POST_EN} = 1. Khác với địa chỉ, cả ba stride luôn bị kiểm, kể cả \sig{LDB} khi \sig{REUSE_B} = 1 \\\hline

Bit và word dự trữ & \sig{FLAGS[31:5]} = 0; \sig{LDC[31:16]} = 0; \sig{QSHIFT[31:5]} = 0; word 12 … 15 = 0 \\\hline

Ràng buộc cờ & \sig{RELU_EN} = 1 kéo theo \sig{POST_EN} = 1 \\\hline

Dùng lại B & \sig{REUSE_B} = 1 chỉ khi đây không phải descriptor đầu chuỗi \textbf{và} K, N bằng của descriptor ngay trước \\\hline

\end{spectable}

Word \sig{TAG} không bị kiểm: mọi giá trị đều hợp lệ và phần cứng không đọc nó.

\paragraph{Chế độ chuỗi}\label{par:5-5}
\begin{itemize}
\item \sig{DESC_PTR} và mọi \sig{NEXT_PTR} được dùng phải căn 64 byte; nếu không, \sig{ERR_CODE} = 1
\item Mỗi descriptor chạy các pha FETCH → kiểm tra → LOAD\_BIAS (nếu \sig{POST_EN}) → LOAD\_A → LOAD\_B (trừ khi \sig{REUSE_B}) → COMPUTE → STORE\_C, với cùng phép toán và cùng bố trí DRAM ở mục~\ref{sec:3}
\item \sig{REUSE_B} = 1: không nạp B; COMPUTE dùng B của descriptor ngay trước, trên đúng bank B của descriptor đó
\item \sig{WAIT_PREV} = 1: các pha LOAD của descriptor chỉ bắt đầu khi STORE\_C của mọi descriptor trước nó đã xong, tức khi \sig{DESC_DONE_CNT} bằng chỉ số của descriptor này. Phần mềm phải đặt cờ này khi descriptor đọc dữ liệu do descriptor trước ghi ra. \sig{WAIT_PREV} chặn các LOAD, không chặn FETCH và kiểm tra
\item Chuỗi kết thúc ở descriptor đầu tiên có \sig{LAST} = 1. DONE = 1 khi STORE\_C của descriptor đó nhận phản hồi B cuối là OKAY
\item \sig{DESC_DONE_CNT} tăng 1 mỗi khi STORE\_C của một descriptor xong OKAY; các descriptor xong đúng theo thứ tự chuỗi
\item Lỗi: gọi f là descriptor có thứ tự nhỏ nhất bị lỗi. Không pha nào của descriptor thứ ≥ f được bắt đầu; các descriptor đứng trước f được chạy xong. Khi mọi khối đã rảnh: ERROR = 1, \sig{ERR_CODE} theo bảng~\ref{tab:2-2} (lỗi FETCH là mã 4), \sig{ERR_DESC_ADDR} = địa chỉ của f, \sig{DESC_DONE_CNT} = số thứ tự của f
\item Khi dùng stride lớn hơn độ dài hàng, DMA cắt burst theo từng hàng: burst không bao giờ trải qua khoảng hở giữa hai hàng, và các byte trong khoảng hở không bị đọc hay ghi
\item Accelerator không ghi vào descriptor; phần mềm không sửa descriptor và vùng dữ liệu của chuỗi đang chạy. Vùng đích của mọi descriptor không được chồng lên vùng descriptor
\item Tham số hậu xử lý của job lấy từ descriptor: \sig{POST_EN}, \sig{RELU_EN}, \sig{BIAS_ADDR}, \sig{QM0}, \sig{QSHIFT}
\end{itemize}

\subsection{Scratchpad ping-pong}\label{sec:5.3}

Mỗi buffer A, B, C, BIAS có hai bank. DMA kênh đọc nạp descriptor d+1 vào một bank trong lúc array
tính descriptor d trên bank kia, và DMA kênh ghi chép kết quả của descriptor d−1 ra DRAM. Thời gian
mỗi descriptor tiến tới $\max(\text{nạp}, \text{tính}, \text{ghi})$ thay vì tổng của chúng.

\begin{spectable}{|L{0.12}|L{0.88}|}{Chọn bank cho descriptor d}{tab:5-6}{\textbf{Buffer} & \textbf{Quy tắc}}
A, C & Đảo bank sau mỗi descriptor \\\hline

B & Đảo sau mỗi lần có LOAD\_B; với \sig{REUSE_B} = 1, dùng lại bank của descriptor trước \\\hline

BIAS & Đảo sau mỗi descriptor có \sig{POST_EN} = 1 \\\hline

\end{spectable}

\subsubsection{Đồng bộ producer và consumer}\label{sec:5.3.1}

Main FSM giữ một \emph{scoreboard} đếm số bên còn sẽ đọc mỗi bank. Đây là cơ chế duy nhất bảo đảm
không có bên nào đọc một bank đang bị ghi, hoặc ghi một bank đang bị đọc.

\begin{spectable}{|L{0.18}|L{0.12}|L{0.7}|}{Scoreboard bank trong Main FSM}{tab:5-7}{\textbf{Bộ đếm} & \textbf{Độ rộng} & \textbf{Ý nghĩa và cập nhật}}
\sig{rcnt[A][b]}, \sig{rcnt[B][b]}, \sig{rcnt[BIAS][b]} & 2 bit & Số job còn sẽ \textbf{đọc} bank b của buffer đó (phía COMPUTE). Tăng ở chu kỳ engine Đọc bắt đầu LOAD\_X của job vào bank b, hoặc ở bước R\_RB với job \sig{REUSE_B} (không tăng lúc gán bank ở CHECK, nếu không LOAD\_X sẽ chờ chính bộ đếm vừa tăng); giảm khi \sig{i_pe_done} của job đó. LOAD\_X vào bank b chỉ được bắt đầu khi \sig{rcnt[X][b]} = 0 \\\hline

\sig{wcnt[C][b]} & 1 bit & Số job còn sẽ \textbf{đọc} bank C b (phía STORE\_C). Tăng khi COMPUTE của job có bank C = b bắt đầu; giảm khi STORE\_C của job đó xong OKAY (\sig{i_wr_done}). COMPUTE ghi bank C b chỉ được bắt đầu khi \sig{wcnt[C][b]} = 0 \\\hline

\end{spectable}

Trọng tài giữa ba engine khi cùng muốn phát lệnh DMA hoặc cùng muốn tiến: ưu tiên cố định
\textbf{Ghi > Tính > Đọc}. Lý do của thứ tự này: engine Ghi giải phóng bank C, engine Tính giải
phóng bank A, B, BIAS; ưu tiên khâu giải phóng tài nguyên trước làm engine Đọc không bị khóa lâu và
giữ đường ống chạy đầy.

\paragraph{Ping-pong}\label{par:5-8}
\begin{itemize}
\item Ba khối thực thi: Đọc (FETCH, kiểm tra, LOAD\_*; kênh đọc DMA), Tính (COMPUTE; array và VPU), Ghi (STORE\_C; kênh ghi DMA). Mỗi khối làm một việc một lúc, theo thứ tự descriptor. Kênh đọc và kênh ghi được phép có burst AXI cùng lúc, mỗi kênh tối đa 1 burst outstanding
\item Khối Đọc làm lần lượt FETCH(d), kiểm tra d, các LOAD của d, rồi mới sang FETCH(d+1)
\item LOAD vào một bank chỉ bắt đầu khi COMPUTE cuối cùng đọc bank đó đã xong
\item COMPUTE(d) bắt đầu khi mọi LOAD của d và COMPUTE(d−1) đã xong, và STORE\_C cuối cùng đọc bank C của d đã xong
\item STORE\_C(d) bắt đầu khi COMPUTE(d) và STORE\_C(d−1) đã xong; do đó kết quả ra DRAM đúng theo thứ tự descriptor
\item \sig{WAIT_PREV} chặn các LOAD, không chặn FETCH và kiểm tra
\item Tại mỗi thời điểm, mỗi bank có tối đa một bên truy cập. Hai bank khác nhau của cùng buffer được truy cập độc lập trong cùng chu kỳ; mỗi bank vẫn chỉ cần SRAM một cổng
\item Các điều kiện thứ tự và quyền truy cập bank nêu trên là điều kiện \textbf{cần}; phần cứng được phép bắt đầu một pha muộn hơn, không bao giờ sớm hơn. Nội dung DRAM sau tác vụ phải giống hệt thực thi tuần tự từng descriptor
\item Soft reset: abort cả hai kênh DMA đang bận; \sig{soft_clr} phát khi cả hai kênh DMA đã rảnh (\sig{i_rd_busy} = 0 và \sig{i_wr_busy} = 0)
\item \sig{PERF_LOAD} + \sig{PERF_COMPUTE} + \sig{PERF_STORE} được phép lớn hơn \sig{PERF_TOTAL}; phần chênh là thời gian chạy chồng
\end{itemize}

\subsubsection{Hành vi biên của ping-pong}\label{sec:5.3.2}

\begin{spectable}{|L{0.26}|L{0.74}|}{Các trường hợp biên}{tab:5-9}{\textbf{Trường hợp} & \textbf{Hành vi}}
Chuỗi một descriptor & Ba engine chạy tuần tự; không có gì để chồng. Bank dùng là bank 0 của mọi buffer \\\hline

Descriptor đầu chuỗi & \sig{rcnt} và \sig{wcnt} đã được xóa ở chu kỳ START được chấp nhận, nên không engine nào bị chặn \\\hline

\sig{REUSE_B} liên tiếp & Bank B không đảo; \sig{rcnt[B][b]} của bank đang dùng tăng thêm cho mỗi job dùng lại. Không có logic dừng riêng cho bộ đếm: theo cấu trúc pipeline, tại mọi thời điểm tối đa 3 job giữ một bank B (một job đang COMPUTE chưa có \sig{i_pe_done}, một job trong slot, một job ở R\_RB). Job ở C\_HAND đã giảm \sig{rcnt} lúc \sig{i_pe_done}; job mới chỉ vào slot khi slot trống, tức sau \sig{i_pe_done} của job trước. Vì vậy \sig{rcnt} 2 bit (tối đa 3) không bao giờ tràn; AS-03 của Main FSM kiểm điều này \\\hline

\sig{POST_EN} xen kẽ & Bank BIAS chỉ đảo ở descriptor có \sig{POST_EN} = 1; descriptor có \sig{POST_EN} = 0 không nạp bias, không đảo bank BIAS và không tăng \sig{rcnt[BIAS]} \\\hline

\sig{WAIT_PREV} = 1 & Engine Đọc đã FETCH và kiểm tra xong nhưng dừng trước LOAD; engine Tính và Ghi vẫn chạy descriptor trước. Ping-pong tạm thu về thực thi tuần tự ở ranh giới đó \\\hline

Lỗi giữa chuỗi & Engine đã bắt đầu pha thì chạy tới hết; không pha mới nào của descriptor ≥ f bắt đầu. Các bộ đếm scoreboard vẫn giảm đúng khi các pha đang chạy kết thúc, nên hệ thống về trạng thái yên và ERROR được đặt \\\hline

Soft reset giữa chuỗi & Không pha mới; abort kênh DMA đang bận; COMPUTE bị hủy ngay. Scoreboard bị xóa bởi \sig{soft_clr}; nội dung scratchpad không xác định \\\hline

\end{spectable}
